\documentclass[10pt]{article}

\usepackage[utf8]{inputenc}

\usepackage[T1]{fontenc}

\usepackage{amsmath}

\usepackage{amsfonts}

\usepackage{amssymb}

\usepackage[version=4]{mhchem}

\usepackage{stmaryrd}

\usepackage{bbold}

\usepackage{CJKutf8}

\usepackage{graphicx}

\usepackage[export]{adjustbox}

\graphicspath{ {./images/} }



\begin{document}

мов бикатегории $\mathfrak{\Re}=(\Re,(\mathbb{F}, \mathfrak{M})$, то $P$ наз. д о п у $\mathrm{c}-$ т п мым пр о ект вв ным обтектом. Напр., в нек-рых многообразиях групп свободные группы этого многообразия являются допустимыми П. о. относительно класса всех сюръективных гомоморфизмов, но не являются П. о., поскольку существуют несюръективные эпиморфизмы.



Дуальным к понятию П. о. является понятие и в ъе к т и в н о г о о б т е к т а. Фундаментальная роль проективных и инъективных объектов была выявлена при построении гомологич. алгебры. В категориях модулей всякий модуль представим в виде фактормодуля проективного модуля. Это свойство позволяөт строить т. н. проективные резольвенты и исследовать различные типы гомологич. размерности.



Лит.: [1] К а р т а н А., Э й л е н б е р г С., Гомологическая алгебра, пер. с англ., М., 1960; [2] М а к л е й н С.,



ПРОЕКТИВНЫЙ ПРЕДЕЛ, о б р а т н Ы й П р ед ө л, - конструкция, возникшая первоначальво в теории множеств и топологии, а затем нашедшая широкое применение во многих разделах математики. Наиболее часто используется П. П. семейства однотипных математич. структур, индексированных әлементами нек-рого предупорядоченного множества. Пусть $I$ - множество, снабженное отношением предпорядка $<$, и каждому әлементу $i \in I$ сопоставлено множество $X_{i}$, а каждой паре $(i, j), i, j \in I$, в к-рой $i<j$, сопоставлено отображение $\varphi_{i j} X_{i} \rightarrow X_{j}$, причем $\varphi_{i i}, i \in I,-$ тождественные отображения и $\varphi_{i j} \varphi_{j k}=\varphi_{i k}$ при $i<j<k$. Множество $X$ наз. п р оективным пре делом се мейс т в а множеств $X_{i}$ и отображений $\varphi_{i j}$, если выполнены следующие условия: a) существует такое семейство отображений $\pi_{i}: X \rightarrow X_{i}$, что $\pi_{i} \varphi_{i j}=\pi_{j}$ для любой пары $i<j ;$ б) для любого семейства отображений $\alpha_{i}: Y \rightarrow$ $\rightarrow X_{i}, \quad i \in I$, произвольного множества $Y$, для к-рого выполнены равенства $\alpha_{i} \varphi_{i j}=\alpha_{j}$ при $i<j$, существует такое однозначно определенное отображение $\alpha: Y \rightarrow X$, что $\alpha_{i}=\alpha \pi_{i}$ для всех $t \in I$. Конструктивно П. П. можно описать следующим образом: рассматривается прямое произведение $\prod_{i \in I} X_{i}$ И в нем выделяөтся подмножество всех функций $f: I \rightarrow \bigcup_{i \in I} X_{i}$, для к-рых выполняются равенства $\varphi_{i j}(f(i))=f(j)$ при $i<j$. Это подмножество является П. П. семейства $X_{i}$. Если все $X_{i}$ снабжены дополнительной однотипной структурой, к-рая переносится на $\prod_{i \in I} X_{i}$, то эта же структура индуцируется и в П. П. Поэтому можно говорить о П. П. групП, модулей, топологич. пространств и т. д.



Естествентым обобщением понятия П. П. является понятие П. п. функтора. Пусть $F: \mathfrak{D} \rightarrow \mathfrak{\Re}-$ одноместный ковариантный функтор из малой категории $\mathfrak{D}$ в произвольную катөгорию \begin{CJK}{UTF8}{mj}凡\end{CJK}. Объект $X \in \mathrm{Ob}$, вместе с морфизмами $\pi_{D}: X \rightarrow F(D), D \in O b \mathfrak{D}$, наз. п р о е ктивным пределом (обратыны пр еде лом, или просто пред өлом) функт о р а $F$, если выполнены следующие условия: а) $\pi_{D} F(\varphi)=$ $=\pi_{D^{\prime}}$ для любого морфизма $\varphi: D \rightarrow D^{\prime} ;$ б) для всякого семейства морфизмов $\alpha_{D}: Y \rightarrow F(D)$, для к-рого $\alpha_{D} F(\varphi)=\alpha_{D}$, при $\varphi: D \rightarrow D^{\prime}$, существует такой единственный морфизм $\alpha: Y \rightarrow X$, что $\alpha_{D}=\alpha \pi_{D}$, для любого $D \in \mathrm{Ob} \mathfrak{D}$. Обозначение: $\lim F=\left(X, \pi_{D}\right)$. Аналогично определяется проективный предел контравариантного функтора.



П р и м ө р ы П. п. 1) Пусть $I$ - дискретная категория. Тогда для произвольного функтора $F: I \rightarrow \Re$ проективный предел функтора $F$ совпадает с прямым произведением семейства объектов $F(i), i \in I$.



\begin{enumerate}

  \setcounter{enumi}{1}

  \item Пусть $\mathfrak{D}$ - категория с двумя объектами $A, B$ и двумя неединичными морфизмами $\alpha, \beta: A \rightarrow B$. Тогда предел любого функтора $F: \mathfrak{D} \rightarrow \mathfrak{R}$ является ядром пары морфизмов $F(\alpha), F(\beta)$. Если в категории $\mathfrak{\Re}$ существуют произведения любых семейств объектов и ядра пар морфизмов, то в $\Re$ существует предел любого функтора $F: \mathfrak{D} \rightarrow \mathbb{\Re}$ из произвольной малой категории $\mathfrak{D}$. $\quad M$. Ш. цаленко. ІРОЕКТИВНЫЙ СПЕКТР к о л ь ц а - схема $X=$ $=\operatorname{Proj}(R)$, сопоставляемая градуированному кольцу $R=\sum_{n=0}^{\infty} R_{n}$. Как множество точек $X$ представляет собою множество однородных простых идеалов $p \subset R$, таких, что $p \not \supset \sum_{n=1}^{\infty} R_{n}$. Топология на $X$ определяется следующим базисом открытых множеств: $X_{f}=\{p \mid f \notin p\}$ для $f \in R_{n}, n>0$. Структурный пучок $\widehat{O} X$ локально окольцованного пространства $X$ задается на базисных открытых множествах так. $\Gamma\left(X_{f}, \sigma_{X}\right)=\left[R_{(f)}\right]_{0}$, т. е. подкольцо элементов степени 0 кольца частных $R_{(f)}$ по мультипликативной системе $\left\{f^{n}\right\}_{n \geqslant 0}$.

\end{enumerate}



Наиболее важным примером П. с. является $p^{n}=$ $=\operatorname{Proj} \mathbb{Z}\left[T_{0}, T_{1}, \ldots, T_{n}\right]$. Множество его $k$-значных точек $P_{k}^{n}$ для любого поля $k$ находится в естествөнном соответствии с множеством точек проективного $n$-мерного пространства над полем $k$.



Если все кольца $R_{m}$ как $R_{0}$-модули натянуты на $\underbrace{R_{1} \otimes \ldots \otimes R_{1}}$, то на Proj $(R)$ определена еще дополни-



тельная структура. А именно, покрытие $\left\{X_{f} \mid f \in R_{1}\right\}$ и



\includegraphics[max width=\textwidth, center]{2023_07_08_ab8af2c5375172152efcg-1(1)}

к-рому отвечает обратимый пучок, обозначаемый через $G(1)$. Через $G(n)$ принято обозначать $n$-ю тензорную степень $\mathcal{O}(1)^{\bigotimes}$ пучка $\mathcal{O}(1)$. Существует канонич.



\includegraphics[max width=\textwidth, center]{2023_07_08_ab8af2c5375172152efcg-1}

метрич. смысл градуировки кольца $R$ (см. [1]). Если, напр., $R=k\left[T_{0}, \ldots, T_{n}\right]$, то $G(1)$ соответствует пучку гиперплоских сечений в $P_{k}^{n}$.



Лum.: [1] М а м ф о p д Д., Лекции о кривых на алгебраической поверхности, пер. с англ., M., 1968; [2] G r o t h e n-



\begin{center}

\includegraphics[max width=\textwidth]{2023_07_08_ab8af2c5375172152efcg-1(2)}

\end{center}



IIPOEKTOP, п р о е к ц и о н ны й о п е р а т о р,линейный оператор $P$ в векторном пространстве $X$ такой, что $P^{2}=P$. $M$. И. Войцеховский.



ПРОЕКЦИОННЫЕ МЕТОДЫ - методы отыскания приближенного решения операторного уравнения в заданном подпространстве, основанные на проектировании уравнения на нек-рое (вообще говоря, другое) подпространство. П. м. являются основой построения различных вычислительных схем решения краевых задач, в том числе метода конечных әлементов и метода коллокации.



Пусть $L$ - оператор, область определения $D(L)$ к-рого лөжит в банаховом пространстве $X$, а область значений $R(L)$ - в банаховом пространстве $Y$. Для решения уравнения



\[

L x=y

\]



проекционным методом выбирают две последовательности подпространств $\left\{X_{n}\right\}$ и $\left\{Y_{n}\right\}$



\[

X_{n} \subset D(L) \subset X, Y_{n} \subset Y, \quad n=1,2, \ldots,

\]



а также проекторы $P_{n}$, проектируюпце $Y$ на $Y_{n}$. Уравнение (1) заменяется приближенным



\[

P_{n} L x_{n}=P_{n} y, x_{n} \in X_{n} .

\]



В случае $X=Y, X_{n}=Y_{n}, n=1 ; 2$, .. П. м. (2) принято называть м е т о д о м $\Gamma$ а л е р к и н а (иногда последний метод трактуется более широко, см. Галеркина метод).



Имеет место т е о р е м а с х о ди м о с т и П. м. для линейных уравнений (для случая конечномерных подпространств $X_{n}$ и $Y_{n}$ ). Пусть $L$ линеен и переводит





\end{document}